\documentclass[11pt,a4paper]{article}
\usepackage[margin=1in]{geometry}
\usepackage{amsmath,amssymb,amsthm}
\usepackage[hyphens]{url}
\usepackage[hidelinks]{hyperref}
\newtheorem{theorem}{Theorem}
\newtheorem{lemma}[theorem]{Lemma}
\newcommand{\Z}{\mathbb Z}
\newcommand{\V}{V}
\title{A counterexample to the fivefold sumset inequality}
\author{Benjamin Guo\\Hraness (\url{hraness.com})\thanks{Developed with AI research agents. The proof is accompanied by two exact integer count formulas and a separate verifier.}}
\date{29 September 2026}
\hypersetup{pdftitle={A counterexample to the fivefold sumset inequality},pdfauthor={Benjamin Guo}}

\begin{document}
\maketitle
\begin{abstract}
We give finite sets of integers $A,B$ satisfying
$|A-B|^5>|A+B|^5|5B|$, where $5B$ is the fivefold sumset of $B$.
The sets are base-$193$ encodings of nonnegative integer vectors in
dimension $128$, with coordinate sums bounded by $160$ and $32$,
respectively. Stars and bars counts the sumsets, and a sum of $33$
binomial terms counts the difference set. The strict comparison uses
integer arithmetic and is independently reproduced by counting both
positive and negative supports.
\end{abstract}
\medskip\noindent\textbf{Keywords.} sumset; difference set; integer simplex; exact certificate.
\smallskip\noindent\textbf{MSC 2020.} 11B75; 05A16.

\section{The inequality and the construction}
For finite nonempty subsets $A,B$ of an abelian group, write
$A\pm B=\{a\pm b:a\in A,\ b\in B\}$ and
$hB=\{b_1+\cdots+b_h:b_1,\ldots,b_h\in B\}$.
Gyarmati, Hennecart and Ruzsa~\cite[equation~(9), p.~177]{ghr} asked for which
integers $h$ the inequality
\begin{equation}\label{eq:historical}
 |A-B|\le |A+B|\,|hB|^{1/h}
\end{equation}
holds for all such sets. Their discussion on p.~182 gives a counterexample
for $h=6$ and leaves $h=4,5$ unresolved. The result below gives an
explicit counterexample to the $h=5$ statement in that text; no priority
or current-status claim is made.

For integers $m\ge1$ and $r\ge0$, define the integer simplex
\[
 \V(m,r)=\left\{x\in\Z_{\ge0}^{m}:
                  \sum_{i=0}^{m-1}x_i\le r\right\}.
\]
Its size is $\binom{m+r}{m}$.
We use unequal radii in the classical lattice-simplex construction,
used by Gyarmati, Hennecart and Ruzsa in their $h=6$
example~\cite[p.~182]{ghr}.

\begin{theorem}\label{thm:main}
Let $\phi:\Z^{128}\to\Z$ be the additive map
\[
 \phi(x)=\sum_{i=0}^{127}193^i x_i.
\]
The finite integer sets
\[
 A=\phi\bigl(\V(128,160)\bigr),\qquad
 B=\phi\bigl(\V(128,32)\bigr)
\]
satisfy
\begin{equation}\label{eq:failure}
 |A-B|^5>|A+B|^5|5B|.
\end{equation}
\end{theorem}

\section{Counting the vector sets}
The count formulas hold for any pair of integer simplices.

\begin{lemma}\label{lem:counts}
For integers $m\ge1$ and $a,b\ge0$, let $X=\V(m,a)$ and $Y=\V(m,b)$. Then
\begin{align}
 |X+Y|&=\binom{m+a+b}{m}, &
 |5Y|&=\binom{m+5b}{m},\label{eq:sumcounts}\\
 |X-Y|&=\sum_{j=0}^{\min(m,b)}
       \binom mj\binom bj\binom{m-j+a}{m-j}.\label{eq:diffcount}
\end{align}
The difference count also equals
\begin{equation}\label{eq:double}
 \sum_{j=0}^{\min(m,b)}\ \sum_{k=0}^{\min(m-j,a)}
 \binom mj\binom{m-j}{k}\binom bj\binom ak.
\end{equation}
\end{lemma}

\begin{proof}
One has $\V(m,a)+\V(m,b)=\V(m,a+b)$. To obtain the reverse
inclusion, allocate up to $a$ units of any vector's coordinate mass
to the first summand; the remaining mass is at most $b$. Iterating
this identity gives $5Y=\V(m,5b)$. Stars and bars proves
\eqref{eq:sumcounts}.

For $z\in\Z^m$, put $(z_+)_i=\max(z_i,0)$ and
$(z_-)_i=\max(-z_i,0)$. Then
\[
 z\in X-Y\quad\Longleftrightarrow\quad
 \sum_i(z_+)_i\le a\ \text{ and }\ \sum_i(z_-)_i\le b.
\]
Necessity follows from any representation $z=x-y$; sufficiency follows
by taking $x=z_+$ and $y=z_-$. If $z$ has exactly $j$ negative
coordinates, their positions can be chosen in $\binom mj$ ways.
Their positive magnitudes have total at most $b$, giving $\binom bj$
choices. The other $m-j$ coordinates are nonnegative with total at
most $a$, giving $\binom{m-j+a}{m-j}$ choices. This proves
\eqref{eq:diffcount}, including the case $j=0$.

Alternatively, choose the $j$ negative positions and then the $k$
positive positions. Their nonzero magnitudes can be filled in
$\binom bj$ and $\binom ak$ ways. Summing over the disjoint supports
proves~\eqref{eq:double}.
\end{proof}

For the parameters of Theorem~\ref{thm:main}, set
\begin{align}
 D&=\sum_{j=0}^{32}
      \binom{128}{j}\binom{32}{j}\binom{288-j}{128-j},\label{eq:D}\\
 S&=\binom{320}{128},\qquad
 F=\binom{288}{128}.\label{eq:SF}
\end{align}
These integers satisfy $D^5>S^5F$. The arithmetic can be checked with
the following standalone Python code, using integer operations only:
\begin{verbatim}
from math import comb
D = sum(comb(128,j)*comb(32,j)*comb(288-j,128-j)
        for j in range(33))
S, F = comb(320,128), comb(288,128)
if D**5 <= S**5*F:
    raise ArithmeticError("fivefold comparison failed")
print("D^5 > S^5 F")
\end{verbatim}
The accompanying verifier also evaluates~\eqref{eq:double} with
factorials instead of the binomial-coefficient routine and checks
the resulting integers against the stored values. Its small-set
controls enumerate the sets themselves.

\section{Encoding as integer sets}
\begin{proof}[Proof of Theorem~\ref{thm:main}]
Write $X=\V(128,160)$ and $Y=\V(128,32)$.
Every coordinate of $X+Y$ lies in $[0,192]$, and every coordinate
of $5Y$ lies in $[0,160]$. Base-$193$ uniqueness therefore makes
$\phi$ injective on both sets.

Every coordinate of $X-Y$ lies in $[-32,160]$. Adding $32$ to
each coordinate gives digits in $[0,192]$ and adds the same integer
$32\sum_{i=0}^{127}193^i$ to every image. Thus $\phi$ is injective
on $X-Y$ as well. Additivity now gives
\[
 |A-B|=D,\qquad |A+B|=S,\qquad |5B|=F.
\]
The integer comparison following~\eqref{eq:SF} proves
\eqref{eq:failure}.
\end{proof}

\begin{thebibliography}{9}
\bibitem{ghr}
K. Gyarmati, F. Hennecart and I. Z. Ruzsa,
\emph{Sums and differences of finite sets},
Functiones et Approximatio Commentarii Mathematici
\textbf{37}(1) (2007), 175--186.
\url{https://doi.org/10.7169/facm/1229618749}.
\end{thebibliography}
\end{document}
